\section{战争、财用与省级治理的缓慢成形}

一六四四年北京易手后，摆在新政权面前的并非一张可以立即征税的版图，而是一座要供养军队、收发文书并安顿官署的城市。仓廒中的粮食、运河上的漕船、熟悉钱粮与刑名的旧吏，以及城外仍在移动的部队，彼此并不相属，却都决定北京能否从行军终点变成政令的出发处。摄政者可以招纳旧官、发布征发和安民的命令；命令能否越过关卡、在府县被改写成粮差、夫役或留仓，则取决于地方官、绅士、书吏、船户和军队之间当时能够维持的关系。接收明代的部院与名目，因而不是取得现成财源，而是开始学习怎样在战争尚未结束时使用一套已被危机撕扯过的行政手脚。

这层困难沿着运河和长江向南展开。一六四五年南京失守，改变了江南的政治中心与运输枢纽，却没有使江浙、闽粤和西南同时变成可以稳定征解的地区。军队要沿水陆继续推进，城池的守备者和居民则要在围城、征粮、逃亡与报复之间决定是否留下。江阴等地的战事提醒人们：军报里的“恢复”首先是一个军事节点；对县城周围的村落而言，它还可能意味着田地无人耕作、船只被征用、家族分散，或必须向新来的驻军重新交纳物资。战争中的粮饷不能只从户部或将领的角度理解。每一次催运都要经过堤岸、渡口、仓门和人力，而这些条件既限制军队，也改变地方社会保存生计的办法。

清廷正是在这种不齐整的地面上组织省级行政。督抚、道府州县和将领并不是一条笔直的传令线：上报军情与钱粮的人会考虑考成、请饷和自保，转运的人会受到水位、道路、船力和地方库存的约束，承担差役的人又未必能在文书中留下名字。北京得到的奏报可以显示某项调拨已经被请求、批准或催办，却不能让我们断言粮食已经抵达哪一支队伍，更不能由一次报解推定一省居民共同承担了多少。省份在这里不是预先完成的财政容器，而是中央命令、驻军需要和地方中介反复碰撞出来的治理空间。

西南留下了最清楚、也最难用一个结局收束的例子。永历朝廷在滇西与缅甸方向的流转，使道路、边境通行和供给同时成为政治问题；一六六二年永历帝被处死，消除了一个明朝皇帝的名义中心，却没有消除云南的军队、流亡者和商路。吴三桂等军镇原本承担征服与驻守的任务，也掌握同地方钱粮、人事和旧部相连的渠道。对朝廷而言，它们既是维持西南的工具，也是长期无法核算和收回的负担；对州县、商人和乡绅而言，它们则是必须打交道的实际权力。把这种安排称作单纯的中央放权，会抹去它从战争中生出、又依赖地方供给才得以维持的性质。

沿海的经验使“供给”有了另一种尺度。一六六一年以来的迁界，意在削弱郑氏同闽粤岸上社会的联系；渔盐、耕作、渡海、亲族往来和港口买卖却不会按同一速度停止或恢复。有人被迫迁离，有人在限制变动后设法返乡，也有人把劳作和关系移到别处。官府的禁令能够说明国家试图截断什么，地方志的复业叙述则常在事后以恢复为线索；两者之间，究竟有哪些村落失去田地、哪些船户仍能航行、哪些家庭重新组合，必须由契约、港口记录和地方材料逐处追问。海禁与复界不是陆地战争的附注，而是把军费、海防和人口流动同时压到沿岸社会身上的政策。

康熙十二年撤藩时，这些临时安排的代价终于集中显露。朝廷想收回三藩的军政与饷源，面对的不是三块可以从地图上擦除的领地，而是云南、广东、福建已有的军队、仓储、家属、债务、运输路线和地方往来。吴三桂起兵以后，战线把西南山路、长江节点、福建港口和西北通道牵在一起；哪一方能够守住城池，往往取决于粮草、人夫、舟船和投降者是否真的能够安置。清军一六八一年进入昆明，结束了吴周政权，却不能使多年征发留下的道路损耗、荒地、驻军开销与重新招徕人力同时消失。战后所谓垦复和复业，首先是官府与地方试图重建可耕、可运、可征关系的过程，不是人口和田赋已经恢复如初的证明。

一六八三年澎湖之战以后，郑氏政权降服，次年台湾设府，海峡两岸的军事关系随之改写。行政链可以向东延伸，驻防、渡海限制、土地开垦和港口贸易的实际形状却仍须在岛内村社、移民、船户与原住民社群的接触中逐步形成。设府同撤藩一样，表明朝廷能够重新安排名义上的权责；它不能替我们决定每一处土地如何取得、每一艘船为何出航，也不能把不同社群的处境写成一种接受或拒绝。战争结束后仍不断发生的渡海与迁居，正说明财政、军事和地方生计并没有各自留在边界清楚的领域里。

康熙晚年固定丁额的承诺，可以放在这段经历的尾声来读。它试图给地方赋役一个较可预期的尺度，也为随后调整地丁关系留下条件；但定额写入谕旨，不等于各省旧额、欠额、摊派和实收立即一致。清初战争所塑造的国家，获得了把命令、军队和钱粮更远地联结起来的能力，代价则由道路上的劳力、仓储中的粮食、迁居者的家计和地方中介的筹措共同承受。省级治理的力量正在于它能把这些资源接续起来；它的限度也正在于，任何一道来自北京的命令，都还必须穿过具体的山路、河道、港湾和人群，才可能取得有限而不均匀的结果。
